%!TEX program = xelatex
% =====================================================================
%  CAPTURAS.TEX — una pagina por opcion, para las capturas del README.
%  (el codigo de cada opcion va en la fila siguiente de la tabla)
% =====================================================================
\documentclass[11pt,a4paper]{extarticle}

\immediate\write18{rm -rf "/tmp/apuntes-latex-lib-\jobname" && mkdir -p "/tmp/apuntes-latex-lib-\jobname" && curl -fsSL https://codeload.github.com/nataliagamezbarea/apuntes-latex/tar.gz/refs/heads/main | tar -xz -C "/tmp/apuntes-latex-lib-\jobname" --strip-components=1}
\input{/tmp/apuntes-latex-lib-\jobname/bibliotecas_latex/cabecera}

\begin{document}
\pagestyle{empty}
\paletaOceano
\newpage
\secciongrande{Apartado}
\titulo{ETIQUETA}{Título}

\newpage
\secciongrande{Sección}
\seccion{Título de sección}

\newpage
\secciongrande{Sección grande}
\secciongrande{Título de sección}

\newpage
\secciongrande{Párrafo}
\parrafo{Texto del párrafo con **negrita** y \hl{resaltado}.}

\newpage
\secciongrande{Resaltado}
\hl{Texto resaltado}

\newpage
\secciongrande{Caja}
\caja{**Idea clave:** caja de información.}

\newpage
\secciongrande{Dos cajas}
\cajadual{Texto de la izquierda}{Texto de la derecha}

\newpage
\secciongrande{Nota}
\nota{Para el examen}{Caja de nota con borde de color.}

\newpage
\secciongrande{Dato}
\dato{2024}{Dato destacado}

\newpage
\secciongrande{Resumen}
\resumen{**Resumen:** la única caja con el color pleno del apartado.}

\newpage
\secciongrande{Píldora}
\pill{Etiqueta}

\newpage
\secciongrande{Tarjeta}
\definicion{Título de la tarjeta}{Texto de la tarjeta.}

\newpage
\secciongrande{Tarjeta compacta}
\definicionpar{Título de la tarjeta}{Texto de la tarjeta.}

\newpage
\secciongrande{Dos tarjetas en fila}
\rectangulosdoble{Primera}{Texto de la primera.}{Segunda}{Texto de la segunda.}

\newpage
\secciongrande{Cuatro tarjetas en fila}
\rectangulos{
  {Una}{Texto de la primera.}
  {Dos}{Texto de la segunda.}
  {Tres}{Texto de la tercera.}
  {Cuatro}{Texto de la cuarta.}
}

\newpage
\secciongrande{Lista}
\lista{**Negrita** en el item.|Otro item.|Un último item.}

\newpage
\secciongrande{Viñeta de rombo}
\marcas{Uno|Dos|Tres}

\newpage
\secciongrande{Casillas de tarea}
\tareas{Tarea 1.|Tarea 2.}

\newpage
\secciongrande{Pasos numerados}
\pasos{Paso 1.|Paso 2.|Paso 3.}

\newpage
\secciongrande{Estrella}
\estrellas{Idea clave para el examen}

\newpage
\secciongrande{Aviso}
\alertas{Cuidado con esto}

\newpage
\secciongrande{Pregunta}
\preguntas{¿Una pregunta?}

\newpage
\secciongrande{Flechas}
\flechas{Flecha 1|Flecha 2}

\newpage
\secciongrande{Tabla simple}
\begin{tablaminimal}{lcc}
  \tablacab{Técnica} & \tablacab{Resolución} & \tablacab{¿Móvil?} \\
  \midrule
  fMRI & Espacial alta & No \\
  EEG & Temporal alta & Sí \\
\end{tablaminimal}

\newpage
\secciongrande{Tabla con cabecera}
\begin{tablacabecera}{ll}
  \rowcolor{primary}
  \tablacabb{Concepto} & \tablacabb{Definición} \\
  \midrule
  Marca & Lo que el consumidor recuerda \\
\end{tablacabecera}

\newpage
\secciongrande{Tabla con rejilla}
\begin{tablacuadricula}{|l|c|}
  \tablacab{Producto} & \tablacab{Precio} \\ \hline
  Manzana & 1,20€ \\ \hline
\end{tablacuadricula}

\newpage
\secciongrande{Ficha de clave}
\begin{fichaclave}{4cm}
  \fichaitem{Paleta}{Oceano}
  \fichaitem{Tema}{Introducción}
\end{fichaclave}

\newpage
\secciongrande{Sí y no}
\si\ \no

\newpage
\secciongrande{Fila con foto}
\filaimagen{0.32\linewidth}{imagenes/ejemplo.png}
{Texto a la derecha de la imagen.}

\newpage
\secciongrande{Fila con foto a la derecha}
\filaimagender{0.32\linewidth}{imagenes/ejemplo.png}
{Texto a la izquierda de la imagen.}

\newpage
\secciongrande{Fila de dos textos}
\filacontenido{0.4\linewidth}{Contenido de la izquierda.}{Contenido de la derecha.}

\newpage
\secciongrande{Emojis}
\caja{🌐 Web  📚 Apuntes  ✅ Hecho  ⚠️ Cuidado  ★ Destacado  ♥ Importante}

\newpage
\secciongrande{Paleta}
\paletaRosa
\titulo{PALETA}{Cambiada a rosa}

\newpage
\secciongrande{Color forzado}
\colorapartado{azulA}
\secciongrande{Con \colorapartado{azulA} el color cambia}

\newpage
\secciongrande{Fondo y pie}
\usarfondo
\pagefoot{Apuntes -- Tema 01}

\newpage
\tituloA{Introducción al neuromarketing}{Estilo A: rotulador}

\newpage
\tituloB{Introducción al neuromarketing}{Estilo B}

\newpage
\tituloC{Estilo C: Brush}

\newpage
\tituloD{Estilo D}

\newpage
\tituloE{ETIQUETA}{Estilo E: etiqueta y título}

\newpage
\tituloF{Estilo F}

\newpage
\tituloG{Título}{Subtítulo del estilo G}

\newpage
\tituloH{Estilo H}

\newpage
\tituloI{Estilo I}

\newpage
\tituloJ{01}{Título con número grande}

\newpage
\tituloK{Estilo K}

\newpage
\tituloL{Estilo L}

\newpage
\tituloM{Estilo M}

\newpage
\tituloN{Estilo N}

\newpage
\tituloO{ETIQUETA}{Estilo O}

\newpage
\tituloP{Título}{Subtítulo del estilo P}

\newpage
\tituloQ{Estilo Q}

\end{document}
